% ============ §30 Connections with derivations (原书 pp.230-245) ============
\section{Connections with derivations}\label{sec:30}\index{Jacobian criterion, section 30}\index{weak Jacobian condition, section 30}\index{Nagata-Zariski-Lipman, section 30}\index{geometrically regular, section 30}

\begin{theorem}[\textbf{Nagata-Zariski-Lipman}]\label{thm:30.1}
Let $(A, m)$ be a complete Noetherian local ring with $\mathbb{Q} \subset A$. Suppose that
$x_1, \ldots, x_r \in m$ and $D_1, \ldots, D_r \in \Der(A)$ are elements satisfying
$\det(D_ix_j) \notin m$. Then
\begin{enumerate}[(i)]
  \item There is a subring $C \subset A$ such that
        \[
          A = C[x_1, \ldots, x_r] \cong C[[X_1, \ldots, X_r]].
        \]
        Therefore $x_1, \ldots, x_r$ are analytically independent over $C$, and $A$ is
        $I$-smooth over $C$, where $I = \sum Ax_i$, and therefore also $m$-smooth over
        $C$.
  \item If $\mathfrak{g} = \sum AD_i$ is a Lie algebra, (that is if $[D_i, D_j] \in \mathfrak{g}$ for all
        $i, j$) then we can take $C$ to be $\{a \in A \mid D_1a = \cdots = D_ra = 0\}$.
\end{enumerate}
\end{theorem}

\begin{proof}
Letting $(c_{ij})$ be the inverse matrix of $(D_ix_j)$, and setting
$D_i' = \sum c_{ij}D_j$, we have $D_i'x_j = \delta_{ij}$, so that we can assume that
$D_ix_j = \delta_{ij}$. Quite generally, for an element $t \in m$ and a derivation
$D \in \Der(A)$, we define a map $E(D, t)\colon A \to A$ by
\[
  E(D, t) = \sum_{n \ge 0} \frac{t^n}{n!}D^n;
\]
by our assumptions, $E(D, t)(a) = \sum (t^n/n!)D^n(a)$ is meaningful, and one sees
easily that $E(D, t)$ is a ring homomorphism. Now set
\[
  E_i = E(D_i, -x_i) \quad \text{and} \quad C_i = \Im(E_i);
\]
then $C_i$ is a subring of $A$, and by computation we see that
\[
  D_i\Big(\sum \frac{(-x_i)^n}{n!}D_i^n(a)\Big) = \sum \frac{(-x_i)^{n-1}}{(n-1)!}D_i^n(a)
      + \sum \frac{(-x_i)^n}{n!}D_i^{n+1}(a) = 0,
\]
so that $C_i \subset \{a \in A \mid D_ia = 0\}$. Conversely, if $D_ia = 0$ then
$E_i(a) = a$, so that $C_i = \{a \in A \mid D_ia = 0\}$. Also, since for any $a \in A$ we have
$E_i(a) \equiv a \mod x_1A$, we see that elements of $A$ can be expanded in power
series in $x_1$ with coefficients in $C_1$, so that $A = C_1[x_1]$. Now
$E_1(x_1) = x_1 - x_1 = 0$, so that $x_1A \subset \Ker E_1$, and conversely, if
$E_1(a) = a - x_1D_1a + \cdots = 0$ then $a \in x_1A$, and therefore $\Ker E_1 = x_1A$.
Also, $D_1 \circ D_1c = 0 \iff E_1(c) = c$, so that $C_1 \cap x_1A = 0$. Now we prove
that $x_1$ is analytically independent over $C_1$; by contradiction, suppose that
\[
  c_\nu x_1^\nu + c_{\nu+1}x_1^{\nu+1} + \cdots = 0 \quad \text{with } c_\nu \in C_1
      \text{ and } c_\nu \ne 0.
\]
Then by Ex.~\hyperref[ex:25.2-ex]{25.2}, since $x_1$ is not a zero-divisor in $A$, we have
$c_\nu \in x_1A$, which is a contradiction. Thus if $0 \ne \varphi(X) \in C_1[X]$ then
$\varphi(x_1) \ne 0$, as required.

If $r > 1$, then write $D_i'$ for the restriction to $C_1$ of $E_1 \circ D_i$; then
$D_i' \in \Der(C_1)$, and $x_j \in C_1$ with $D_i'x_j = \delta_{ij}$ for $2 \le i, j \le r$, so that by
induction we have
\[
  C_1 = C[x_2, \ldots, x_r] \cong C[[X_2, \ldots, X_r]].
\]
(i) follows from this.

If $\mathfrak{g}$ is a Lie algebra, then we first arrange as before that $D_ix_j = \delta_{ij}$,
and then set $[D_i, D_j] = \sum_\nu a_{ij\nu}D_\nu$ with $a_{ij\nu} \in A$; then
$[D_i, D_j]x_\nu = D_i(\delta_{j\nu}) - D_j(\delta_{i\nu}) = 0$, so that $a_{ij\nu} = 0$, hence
$[D_i, D_j] = 0$, and $D_1(D_i(C_1)) = D_i(D_1(C_1)) = D_i(0) = 0$. Therefore
$D_i(C_1) \subset C_1$ for $i > 1$, and then in the above notation $D_i' = D_i$ for $i > 1$.
Thus by induction we have $C = \{a \in A \mid D_1a = \cdots = D_ra = 0\}$.
\end{proof}

\begin{corollary*}
Let $(A, m)$ be a reduced $n$-dimensional local ring containing $\mathbb{Q}$, and suppose
that the completion $\hat{A}$ of $A$ is also reduced. If there exist $n$ elements
$D_1, \ldots, D_n \in \Der(A)$ and $x_1, \ldots, x_n \in m$ such that
$\det(D_ix_j) \notin m$ then $A$ is a regular local ring and $x_1, \ldots, x_n$ is a regular
system of parameters of $A$. Suppose in addition that $\mathfrak{g} = \sum AD_i$ is a Lie
algebra; then $C = \{a \in A \mid D_1a = \cdots = D_na = 0\}$ is a coefficient field of
$\hat{A}$.
\end{corollary*}

\begin{proof}
Consider $\hat{A}$, with each of the $D_i$ extended to $\hat{A}$. If
$[D_i, D_j] = \sum_\nu a_{ij\nu}D_\nu$ holds in $A$ then it also holds in $\hat{A}$, so that if
$\mathfrak{g}$ is a Lie algebra, so is $\hat{\mathfrak{g}} = \sum \hat{A}D_i$. By the theorem,
$\hat{A} = C[[x_1, \ldots, x_n]]$, and $C$ is isomorphic to $\hat{A}/\sum x_i\hat{A}$, so is a
zero-dimensional local ring. Now by assumption $C$ is also reduced, so that $C$
must be a field. Therefore $\hat{A}$ is a formal power series ring over a field, and
hence is regular, so that $A$ is also regular. The other assertion is also clear.
\end{proof}

\begin{remark}
If we view this corollary as a criterion for regularity, then the condition that
$\hat{A}$ should be reduced is rather a nuisance; however, as we will see later, for
a very wide variety of local rings, we have $A$ is reduced $\iff \hat{A}$ is reduced.
This is the case (corollary of Theorem~\hyperref[thm:32.6]{32.6}) if $A$ is a localisation of a ring
$B$ which is finitely generated over a field (such an $A$ is said to be
\emph{essentially of finite type} over $K$). Note also that if we start off with a
regular local ring $A$, then the corollary gives a concrete method of
constructing a coefficient field of $\hat{A}$.
\end{remark}

Next we consider rings which are finitely generated over a field, which are
important in algebraic geometry.

\begin{theorem}\label{thm:30.2}
Let $k$ be a field, and $A = k[x_1, \ldots, x_n]$ a finitely generated ring over $k$. If
$A_p$ is 0-smooth over $k$ for every $p \in \mathrm{m}$-$\Spec A$, then $A$ is 0-smooth over
$k$.
\end{theorem}

\begin{proof}
Write $k[X]$ for $k[X_1, \ldots, X_n]$, and let $I = \{f(X) \in k[X] \mid f(x) = 0\}$, so
that $A = k[X]/I$. Suppose that $I = (f_1, \ldots, f_s)$. Consider a commutative diagram
\[
\begin{tikzcd}
A \arrow[r, "\bar{\varphi}"] & C/N \\
k \arrow[u] \arrow[r] & C \arrow[u]
\end{tikzcd}
\]
where $C$ is a ring, and $N \subset C$ is an ideal satisfying $N^2 = 0$. To lift
$\bar{\varphi}$ to $A \to C$, we first of all choose $u_i \in C$ such that
$\bar{\varphi}(x_i) = \bar{u}_i$. If $f \in I$ then $f(u_1, \ldots, u_n) \in N$. Now if we can choose
$y_i \in N$ for $1 \le i \le n$ such that $f_j(u + y) = 0$ for all $j$, then the
homomorphism $A \to C$ defined by $X_i \mapsto u_i + y_i$ is a lifting of
$\bar{\varphi}$. We have
\[
  f_j(u + y) = f_j(u) + \sum_{i=1}^{n} (\partial f_j/\partial X_i)(u) \cdot y_i,
\]
so that we are looking for solutions in $N$ of the system of linear equations in
$y_1, \ldots, y_n$:
\[
  (*) \quad f_j(u) + \sum_{i=1}^{n} (\partial f_j/\partial X_i)(u) \cdot y_i = 0
      \quad \text{for } j = 1, \ldots, s.
\]
For each maximal ideal $p$, the local ring $A_p$ is 0-smooth over $k$, so that if we
set $S = \bar{\varphi}(A - p)$ and $\tilde{S} = \varphi^{-1}(S)$ then in the diagram
\[
\begin{tikzcd}
A_p \arrow[r, "\bar{\varphi}_p"] & C_{\tilde{S}}/N_{\tilde{S}} = C_S/N_S \\
k \arrow[u] \arrow[r] & C_{\tilde{S}} \arrow[u]
\end{tikzcd}
\]
there exists a $y^p\colon A_p \to C_{\tilde{S}}$ lifting $\bar{\varphi}_p$. From this we see that (*), as
a system of equations in $N_{\tilde{S}}$, has a solution in $N_{\tilde{S}}$. If we view $N$ as a
$C/N$-module then $N_{\tilde{S}} = N_S$. Thus the theorem reduces to the following
lemma.
\end{proof}

\begin{lemma}\label{lem:30.1}
Let $A$ and $B$ be rings, $\varphi\colon A \to B$ a ring homomorphism, and $N$ a
$B$-module; suppose that $b_{ij} \in B$ and $\beta_i \in N$. If the system of linear
equations
\[
  \sum_j b_{ij}y_j = \beta_i \quad (\text{for } i = 1, \ldots, s)
\]
has a solution in $N_{A-p}$ for every $p \in \mathrm{m}$-$\Spec A$, then it has a solution in $N$.
\end{lemma}

\begin{proof}
The assumption that there is a solution in $N_{A-p}$ means that there exist
$q_{jp} \in N$ for $1 \le j \le n$, and $t_p \in A - p$ such that
\[
  \sum_j b_{ij}q_{jp} - \varphi(t_p)\beta_i = 0 \quad \text{for } 1 \le i \le s.
\]
Now since $\sum_p t_pA = A$, there is a finite set $\{p_1, \ldots, p_\ell\} \subset
\mathrm{m}$-$\Spec A$ and $a_\nu \in A$ such that
\[
  \sum_{\nu=1}^{\ell} a_\nu t_{p_\nu} = 1.
\]
Hence if we set
\[
  q_j = \sum_\nu \varphi(a_\nu)q_{jp_\nu},
\]
we get $\sum_j b_{ij}q_j = \beta_i$ for $1 \le i \le s$.
\end{proof}

\begin{theorem}\label{thm:30.3}
Let $k$ be a field, $S = k[X_1, \ldots, X_n]$, and $I$, $P$ ideals of $S$ such that
$I \subset P \in \Spec S$. Set
\[
  S_P = R, \quad \rad(R) = PR = M, \quad \text{and} \quad R/IR = A, \quad \rad(A) = m,
\]
\[
  R/M = A/m = K,
\]
and suppose that $\htop IR = r$ and $I = (f_1(X), \ldots, f_t(X))$. Then the following
conditions are equivalent.
\begin{enumerate}[(1)]
  \item $\rank(\partial f_i/\partial X_j \bmod P)_{i=1, \ldots, t;\, j=1, \ldots, n} = r$;
  \item $A$ is 0-smooth over $k$;
  \item $A$ is $m$-smooth over $k$;
  \item $\Omega_{A/k}$ is a free $A$-module of rank $n - r$;
  \item $A$ is an integral domain, its field of fractions is separable over $k$, and
        $\Omega_{A/k}$ is a free $A$-module.
\end{enumerate}
\end{theorem}

\begin{proof}
Note that $R$ is a regular local ring.

(1) $\Rightarrow$ (2) By assumption, for a suitable choice of $r$ elements
$D_1, \ldots, D_r$ from $\partial/\partial X_1, \ldots, \partial/\partial X_n$, and of $r$ elements
$g_1, \ldots, g_r$ from $f_1, \ldots, f_t$, we have $\det(D_ig_j) \notin M$. We observe that
taking $R \ni f$ into $(\partial f/\partial X_1 \bmod M, \ldots, \partial f/\partial X_n \bmod M)
\in K^n$ induces a linear map $M/M^2 \to K^r$, so that the images of $g_1, \ldots, g_r$ in
$M/M^2$ are linearly independent. Therefore $\sum g_iR$ is a height $r$ prime ideal
contained in $IR$, and hence $\sum g_iR = IR$. Given a commutative diagram
\[
\begin{tikzcd}
A \arrow[r, "\bar{\varphi}"] & C/N \\
k \arrow[u] \arrow[r] & C \arrow[u]
\end{tikzcd}
\]
with $N^2 = 0$, write $\bar{X}_i \in A$ for the image of $X_i$, and choose $u_i \in C$
such that $\varphi(u_i) = \bar{\varphi}(\bar{X}_i) \in C/N$. Then there is a homomorphism
$R \to C$ defined by $X_i \mapsto u_i$, and this induces a lifting of $\bar{\varphi}$ to $C$ if and
only if $g_i(u) = 0$ for $1 \le i \le r$. Therefore, as in the proof of the previous
theorem, we need only solve the system of equations in unknowns
$y_1, \ldots, y_n \in N$:
\[
  (*) \quad g_i(u) + \sum_{j=1}^{n} (\partial g_i/\partial X_j)(u) \cdot y_j = 0
      \quad \text{for } 1 \le i \le r.
\]
However, we view $N$ as a $C/N$-module, then via $\bar{\varphi}$ as an $A$-module, so that we
can replace the coefficients $(\partial g_i/\partial X_j)(u)$ of this system by
$(\partial g_i/\partial X_j)(\bar{X}) \in A$. Now by assumption, one $r \times r$ minor of this
$r \times n$ matrix of coefficients is a unit of $A$, so that (*) can always be solved.

(2) $\Rightarrow$ (3) is trivial.

(3) $\Rightarrow$ (1) By \S\ref{sec:28}, Lemma~\hyperref[lem:28.1]{1}, $A$ is regular, so that $IR$ is generated by
$r$ elements forming a subset of a regular system of parameters of $R$, and the
image of the natural map $IR \to M \to M/M^2$ is an $r$-dimensional $K$-vector space.
Let $k_1 \subset k$ be the prime subfield; then by
Theorem~\hyperref[thm:26.9]{26.9} $K = R/M$ is 0-smooth over $k_1$, so that by
Theorem~\hyperref[thm:25.2]{25.2}, the sequence
\[
  0 \to M/M^2 \to \Omega_{R/k_1} \otimes_R K \to \Omega_{K/k_1} \to 0
\]
is exact. We can write $\Omega_{S/k_1} = (\Omega_{k/k_1} \otimes S) \oplus F$, where $F$ is the free
$S$-module with basis $dX_1, \ldots, dX_n$ (for example, by
Theorem~\hyperref[thm:25.1]{25.1}), and localising we get
\[
  \Omega_{R/k_1} = (\Omega_{k/k_1} \otimes R) \oplus F \otimes_S R.
\]
Hence $\Omega_{R/k_1} \otimes_R K = (\Omega_{k/k_1} \otimes_k K) \oplus (F \otimes K)$. However, from
\[
  0 \to I/I^2 \to \Omega_{S/k_1} \otimes_S A \to \Omega_{A/k_1} \to 0
\]
we get the exact sequence $(I/I^2) \otimes_A K \to \Omega_{R/k_1} \otimes_R K \to
\Omega_{A/k_1} \otimes_A K \to 0$, and if $A$ is $m$-smooth over $k$ then by the corollary to
Theorem~\hyperref[thm:28.6]{28.6}, $\Omega_{k/k_1} \otimes K \to \Omega_{A/k_1} \otimes K$ is injective, so that
$V = \Im((I/I^2) \otimes K \to \Omega_{R/k_1} \otimes K)$ maps isomorphically to its
projection $W \subset F \otimes K$ in the second factor of the decomposition
$\Omega_{R/k_1} \otimes K = (\Omega_{k/k_1} \otimes K) \oplus (F \otimes K)$. Now factor
$I/I^2 \otimes K \to \Omega_{R/k_1} \otimes K$ as the composite
$I/I^2 \otimes K \to M/M^2 \to \Omega_{R/k_1} \otimes K$; as we have seen above, the first arrow
has rank $r$, and the second is injective, so that $\rank V = r$. Now
$F \otimes K = KdX_1 + \cdots + KdX_n$, and if we write
$(\partial f_i/\partial X_j) \bmod P = a_{ij}$ then $W$ is spanned by $\sum_j a_{ij}dX_j$ for
$1 \le i \le t$. Therefore $\rank(a_{ij}) = r$, and this proves (1).

(2) $\Rightarrow$ (4) By Theorem~\hyperref[thm:25.2]{25.2},
\[
  0 \to IR/I^2R \to \Omega_{R/k} \otimes_R A \to \Omega_{A/k} \to 0
\]
is a split exact sequence, and since $\Omega_{R/k} \otimes_R A$ is a free $A$-module with basis
$dX_1, \ldots, dX_n$, the $A$-module $\Omega_{A/k}$ is projective; but $A$ is a local ring, so
that $\Omega_{A/k}$ is free. Also $A$ is a regular local ring, therefore an integral
domain, and if $L$ is its field of fractions then $L$ is 0-smooth over $A$, hence also
0-smooth over $k$, so that $L$ is separable over $k$.

(4) $\Rightarrow$ (5) By Theorem~\hyperref[thm:26.2]{26.2}, the field of fractions $L$ of $A$ is
separably generated over $k$, so that a separating transcendence basis of $L$ over
$k$ is a differential basis, and $\rank \Omega_{L/k} = \trdeg_k L = n - r$; but
$\Omega_{L/k} = \Omega_{A/k} \otimes_A L$ and hence $\rank \Omega_{A/k} = \rank \Omega_{L/k} = n - r$.

(5) $\Rightarrow$ (1) In the exact sequence
\[
  IR/I^2R \to \Omega_{R/k} \otimes_R A \to \Omega_{A/k} \to 0,
\]
set $E = \Im(IR/I^2R \to \Omega_{R/k} \otimes_R A)$. Then
$\Omega_{R/k} \otimes_R A \cong E \oplus A^{n-r}$, so that $E \cong A^r$, and therefore
$E \otimes K \cong K^r$. This gives (1).
\end{proof}

\begin{remark}
In the above proof, the equivalence of (1), (2), (4) and (5) was comparatively
easy. The proof of (3) $\Rightarrow$ (1) used the corollary to Theorem~\hyperref[thm:28.6]{28.6}, and so
is not very elementary.
\end{remark}

\begin{remark}
If $k$ is a perfect field, or more generally if the residue field $K$ is separable
over $k$, then $m$-smoothness is equivalent to $A$ being a regular local ring, so
that Theorem~\hyperref[thm:30.3]{3} gives a criterion for regularity. In the case of an imperfect
field $k$, if $A$ is 0-smooth over $k$ then so is the field of fractions $L$ of $A$, but
the residue field $K$ is not necessarily separable over $k$; for example,
$A = k[X]_{(X^p - a)}$ with $a \in k - k^p$. For the case of an imperfect field $k$, we
give a regularity criterion for $A$ in Theorem~\hyperref[thm:30.5]{5}.
\end{remark}

Quite generally, let $A$ be a ring and $P$ a prime ideal of $A$, and let
$D_1, \ldots, D_s \in \Der(A)$ and $f_1, \ldots, f_t \in A$; then we write
\[
  J(f_1, \ldots, f_t; D_1, \ldots, D_s)(P) = (D_if_j \bmod P).
\]
This is an $s \times t$ matrix with entries in the integral domain $A/P$.

\begin{theorem}\label{thm:30.4}
Let $R$ be a regular ring, $P \in \Spec R$, and let $I \subset P$ be an ideal of $R$;
suppose that $\htop IR_P = r$.
\begin{enumerate}[(i)]
  \item for any $D_1, \ldots, D_s \in \Der(R)$ and $f_1, \ldots, f_t \in I$ we have
        \[
          \rank J(f_1, \ldots, f_t; D_1, \ldots, D_s)(P) \le r;
        \]
  \item if $D_1, \ldots, D_r \in \Der(R)$ and $f_1, \ldots, f_r \in I$ are such that
        $\det(D_if_j) \notin P$ then $IR_P = (f_1, \ldots, f_r)R_P$ and $R_P/IR_P$ is regular.
\end{enumerate}
\end{theorem}

\begin{proof}
(i) If $Q$ is a prime ideal of $R$ with $I \subset Q \subset P$ and $\htop Q = r$ then
\[
  \rank J(f_1, \ldots; D_1, \ldots)(P) \le \rank J(f_1, \ldots; D_1, \ldots)(Q),
\]
and if we set $QR_Q = \mathfrak{n}$ then $R_Q$ is an $r$-dimensional regular local ring, so
that $m$ is generated by $r$ elements, $m = (g_1, \ldots, g_r)$. Working in $R_Q$, we can
write
\[
  f_j = \sum_{\nu=1}^{r} g_\nu a_{\nu j} \quad \text{with } a_{\nu j} \in R_Q,
\]
so that
\[
  D_if_j \equiv \sum_{\nu=1}^{r} (D_ig_\nu) \cdot a_{\nu j} \bmod Q,
\]
and therefore
\[
  \rank J(f_1, \ldots, f_t; D_1, \ldots, D_s)(P) \le
  \rank J(g_1, \ldots, g_r; D_1, \ldots, D_s)(Q) \le r.
\]
(ii) Set $M = PR_P$; then if $\det(D_if_j) \notin M$ one sees easily that the images in
$M/M^2$ of $f_1, \ldots, f_r$ are linearly independent over $R_P/M = \kappa(M)$, so that
$\sum f_iR_P$ is a prime ideal of height $r$, and therefore coincides with $IR_P$.
Also, $R_P/IR_P$ is regular.
\end{proof}

\begin{theorem}[\textbf{Zariski}]\label{thm:30.5}
Let $k$ be a field of characteristic $p$, $S = k[X_1, \ldots, X_n]$, and $I$ and $P$
ideals of $S$ with $I \subset P \in \Spec S$. Set $S_P = R$ and $R/IR = A$, and suppose
that $\htop IR = r$ and $I = (f_1, \ldots, f_t)$. Then the following three conditions are
equivalent.
\begin{enumerate}[(1)]
  \item $A$ is a regular local ring.
  \item For any $p$-basis $\{u_\gamma\}_{\gamma \in \Gamma}$ of $k$, define $D_\gamma \in \Der(S)$ by
        $D_\gamma(u_{\gamma'}) = \delta_{\gamma\gamma'}$ (the Kronecker $\delta$) and
        $D_\gamma X_i = 0$; then there are a finite number of elements
        $\alpha, \beta, \ldots, \gamma \in \Gamma$ such that
        \[
          \rank J(f_1, \ldots, f_t; D_\alpha, D_\beta, \ldots, D_\gamma,
              \partial/\partial X_1, \ldots, \partial/\partial X_n)(P) = r.
        \]
  \item There exists a subfield $k' \subset k$ with the following properties:
        $k^p \subset k' \subset k$, $[k : k'] < \infty$, and $\Omega_{A/k'}$ is a free $A$-module, with
        \[
          \rank \Omega_{A/k'} = n - r + \rank \Omega_{k/k'}.
        \]
\end{enumerate}
\end{theorem}

\begin{proof}
(2) $\Rightarrow$ (1) comes from Theorem~\hyperref[thm:30.4]{4}, (ii).

(1) $\Rightarrow$ (2), (3) If $A$ is regular then $IR$ is generated by $r$ elements, and
these form an $R$-sequence, so that $IR/I^2R$ is a free $A$-module of rank $r$. Set
$M = PR$, $m = M/IR$ and $K = R/M = A/m$; then the image of $IR \to M/M^2$ is an
$r$-dimensional $K$-vector space, so that the natural map $(IR/I^2R) \otimes_A K \to
M/M^2$ is injective. From the exact sequence
\[
  IR/I^2R \to \Omega_{R/k} \otimes_R A \to \Omega_{A/k} \to 0
\]
we get the exact sequence
\[
  (*) \quad (IR/I^2R) \otimes_A K \to \Omega_{R/k} \otimes_R K \to \Omega_{A/k} \otimes_A K \to 0.
\]
Now by Theorem~\hyperref[thm:25.2]{25.2},
\[
  0 \to M/M^2 \to \Omega_{R/k} \otimes_R K \to \Omega_{K/k} \to 0
\]
is an exact sequence. The first arrow in (*) is the composite
$(IR/I^2R) \to M/M^2 \to \Omega_{R/k} \otimes K$, and so is injective. Thus
\[
  (**) \quad 0 \to (IR/I^2R) \otimes K \to \Omega_{R/k} \otimes K \to \Omega_{A/k} \otimes K \to 0
\]
is exact. Let $\{du_\sigma\}_{\sigma \in \Sigma}$ be a basis of $\Omega_k$ over $k$; then $\Omega_S$ is the free
$S$-module with basis $\{du_\sigma\}_{\sigma \in \Sigma} \cup \{dX_1, \ldots, dX_n\}$, and
$\Omega_R = \Omega_S \otimes_S R$, $\Omega_R \otimes_R K = \Omega_k \otimes_k K \oplus
(KdX_1 \oplus \cdots \oplus KdX_n)$. Now reorder $f_1, \ldots, f_t$ so that $f_1, \ldots, f_r$ are
generators of $IR$; then $df_1, \ldots, df_r \in \Omega_R \otimes K$ can be expressed using
$dX_1, \ldots, dX_n$ together with finitely many elements $du_{\tau_1}, \ldots, du_{\tau_s}$, and this
gives (2). Now if we let $k'$ be the field obtained by adjoining to $k^p$ all the
elements $u_\sigma$ for $\sigma \in \Sigma$ other than the $du_{\tau_1}, \ldots, du_{\tau_s}$ just used, then
from (**) we see that
\[
  0 \to (IR/I^2R) \otimes K \to \Omega_{R/k'} \otimes K \to \Omega_{A/k'} \otimes K \to 0
\]
is also exact. In the exact sequence
\[
  (\dagger) \quad IR/I^2R \to \Omega_{R/k'} \otimes_R A \to \Omega_{A/k'} \to 0,
\]
the middle term is the free $A$-module with basis
$du_{\tau_1}, \ldots, du_{\tau_s}, dX_1, \ldots, dX_n$; now the generators $f_1, \ldots, f_r$ of $IR$ map
to $df_1, \ldots, df_r \in \Omega_{R/k'} \otimes_R A$, and since these are linearly independent over
$K$, they base a direct summand of $\Omega_{R/k'} \otimes_R A$. Thus $\Omega_{A/k'}$ is also a
free $A$-module, and $\rank \Omega_{A/k'} + r = \rank \Omega_{R/k'} = \rank \Omega_{k/k'} + n$.

(3) $\Rightarrow$ (1) In the exact sequence ($\dagger$) the two terms
$\Omega_{R/k'} \otimes_R A$ and $\Omega_{A/k'}$ are both free modules, and the difference between
their ranks is $r$, so that $IR/I^2R$ maps onto a direct summand of
$\Omega_{R/k'} \otimes_R A$, which is a free module of rank $r$. Thus we can choose
$f_1, \ldots, f_r \in IR$ such that $df_1 \otimes 1, \ldots, df_r \otimes 1$ base this direct summand,
and then by NAK we see that $Rdf_1 + \cdots + Rdf_r$ is a direct summand of
$\Omega_{R/k'}$. Hence there exist $D_1, \ldots, D_r \in \Der_{k'}(R)$ such that
$\det(D_if_j) \notin M$. Therefore, by Theorem~\hyperref[thm:30.4]{4}, $A$ is regular.
\end{proof}

\begin{corollary*}
Let $k$ be a field and $S = k[X_1, \ldots, X_n]$; let $I$ be an ideal of $S$, and set
$B = S/I$. Define $U = \{p \in \Spec B \mid B_p \text{ is 0-smooth over } k\}$ and
$\Reg(B) = \{p \in \Spec B \mid B_p \text{ is regular}\}$. Then both $U$ and $\Reg(B)$ are open
subsets of $\Spec B$.
\end{corollary*}

\begin{proof}
Set $V' = \Spec B$, and let $V_1, \ldots, V_s$ be the irreducible components of $V'$. To
say that $P \in V_i \cap V_j$ for $i \neq j$ means just that $B_p$ has at least two minimal
prime ideals, and such points cannot belong to $\Reg(B)$, (nor a fortiori to $U$).
Thus first of all we can throw out the closed subset $W = \bigcup_{i \neq j}(V_i \cap V_j)$,
and therefore we need only prove that $V_i \cap U$ and $V_i \cap \Reg(B)$ are open in
$V_i$ for each $i$; we fix $i$, and set $\dim V_i = n - r$. Then by
Theorems~\hyperref[thm:30.3]{3} and \hyperref[thm:30.4]{4}, if we let $\Delta_1, \ldots, \Delta_s$ denote the images in
$B$ of the $r \times r$ minors of the Jacobian matrix $(\partial f_i/\partial X_j)$, where
$I = (f_1, \ldots, f_t)$, then $V_i - U$ is the intersection of $V_i$ with the closed subset of
$V'$ defined by the ideal $(\Delta_1, \ldots, \Delta_s)B$ of $B$, and is thus closed in $V_i$.
Using Theorem~\hyperref[thm:30.5]{5}, we can argue similarly for $\Reg(B)$; the only difference
is that, instead of one Jacobian matrix, we have to consider the closed subsets
of $V'$ defined by the ideal of $B$ generated by all the $r \times r$ minors of the
infinitely many Jacobian matrices $J(f_1, \ldots, f_t; D_\alpha, D_\beta, \ldots, D_\gamma,
\partial/\partial X_1, \ldots, \partial/\partial X_n)$, where $\{D_\alpha, \ldots, D_\gamma\}$ runs through all finite
subsets of the set of derivations $\{D_\gamma\}_{\gamma \in \Gamma}$ appearing in
Theorem~\hyperref[thm:30.5]{5}, (2).
\end{proof}

Theorems~\hyperref[thm:30.3]{3} and \hyperref[thm:30.5]{5} contain the result known as the \emph{Jacobian
criterion for regularity} in polynomial rings. There is some purpose in trying to
extend this to more general rings. In cases when the module of differentials is
not finitely generated, then the above method cannot be used as it stands, so we
approach the problem using modules of derivations.

Quite generally, let $A$ be an integral domain with field of fractions $L$; then
for an $A$-module $M$, we write $\rank_A M$ for the dimension over $L$ of the vector
space $M \otimes_A L$.

\begin{theorem}[\textbf{M.~Nomura}]\label{thm:30.6}
Let $(R, m)$ be an equicharacteristic $n$-dimensional regular local ring, and $R^*$
the completion of $R$; suppose that $k$ is a quasi-coefficient field of $R$ and $K$ a
coefficient field of $R^*$ containing $k$. Let $x_1, \ldots, x_n$ be a regular system of
parameters of $R$.
\begin{enumerate}[(i)]
  \item $R^* = K[[x_1, \ldots, x_n]]$, and if we write $\partial/\partial x_i$ for the partial
        derivatives in this representation, then $\Der_k(R^*) = \Der_k(R^*)$ is the free
        $R^*$-module with basis $\partial/\partial x_i$ for $1 \le i \le n$.
  \item The following conditions are equivalent:
        \begin{enumerate}[(1)]
          \item $\partial/\partial x_i$ maps $R$ into $R$ for $1 \le i \le n$, so that they can be
                considered as elements of $\Der_k(R)$;
          \item there exist $D_1, \ldots, D_n \in \Der_k(R)$ and $a_1, \ldots, a_n \in R$ such
                that $D_ia_j = \delta_{ij}$;
          \item there exist $D_1, \ldots, D_n \in \Der_k(R)$ and $a_1, \ldots, a_n \in R$ such
                that $\det(D_ia_j) \notin m$;
          \item $\Der_k(R)$ is a free $R$-module of rank $n$;
          \item $\rank \Der_k(R) = n$.
        \end{enumerate}
\end{enumerate}
\end{theorem}

\begin{proof}
(i) Since $K$ is 0-\'etale over $k$, any derivation of $R^*$ which vanishes on $k$ also
vanishes on $K$. If $D \in \Der_k(R^*)$, set $Dx_i = y_i$; then for any
$f \in R^* = K[[x_1, \ldots, x_n]]$ we have $D(f) = \sum_{i=1}^{n}(\partial f/\partial x_i) \cdot y_i$,
and hence $D = \sum y_i \partial/\partial x_i$. Conversely, for any given $y_i$ we can
construct a derivation by this formula, so that $\Der_k(R^*)$ is the free $R^*$-module
with basis $\partial/\partial x_1, \ldots, \partial/\partial x_n$.

(ii) (1) $\Rightarrow$ (2) $\Rightarrow$ (3) and (4) $\Rightarrow$ (5) are trivial. If (3) holds then
$D_1, \ldots, D_n$ are linearly independent over $R$ and over $R^*$, so that by (i), any
$D \in \Der_k(R)$ can be written as a combination $D = \sum c_iD_i$ of the $D_i$ with
coefficients in the field of fractions of $R^*$, but since $Da_j = \sum c_iD_i(a_j)$ for
$j = 1, \ldots, n$, we have $c_i \in R$; therefore $D_1, \ldots, D_n$ form a basis of
$\Der_k(R)$, which proves (4).

(5) $\Rightarrow$ (1) If $D_1, \ldots, D_n$ are linearly independent over $R$ then there exist
$a_1, \ldots, a_n \in R$ such that $\det(D_ia_j) \ne 0$. Therefore $D_1, \ldots, D_n$ are also
linearly independent over $R^*$. Thus writing $L^*$ for the field of fractions of
$R^*$ we can write $\partial/\partial x_i = \sum_j c_{ij}D_j$, with $c_{ij} \in L^*$. From this we
get $\delta_{ij} = \sum_\nu c_{i\nu}D_\nu x_j$, so that the matrix $(c_{ij})$ is the inverse of
$(D_jx_i)$, and $c_{ij} \in L$, where $L$ is the field of fractions of $R$; therefore
$(\partial/\partial x_i)(R) \subset L \cap R^* = R$.
\end{proof}

\begin{lemma}\label{lem:30.2}
Let $R$ be a regular ring and $P \in \Spec R$ with $\htop P = r$; then the following
two conditions are equivalent:
\begin{enumerate}[(1)]
  \item there exist $D_1, \ldots, D_r \in \Der(R)$ and $f_1, \ldots, f_r \in P$ such that
        $\det(D_if_j) \notin P$;
  \item for all $Q \in \Spec R$ with $\htop Q = s \le r$ such that $Q \subset P$ and
        $R_P/QR_P$ is regular, there exist $D_1, \ldots, D_s \in \Der(R)$ and
        $g_1, \ldots, g_s \in Q$ such that $\det(D_ig_j) \notin P$.
\end{enumerate}
\end{lemma}

\begin{proof}
(2) contains (1) as the special case $P = Q$; conversely, suppose that (1) holds.
By Theorem~\hyperref[thm:30.4]{4} we see that $f_1, \ldots, f_r$ is a regular system of parameters
of $R_P$. If $R_P/QR_P$ is regular then we can take $g_1, \ldots, g_s \in Q$ forming a
minimal basis of $QR_P$, and then take $g_{s+1}, \ldots, g_r \in P$ such that
$g_1, \ldots, g_r$ is a regular system of parameters of $R_P$. Then from
$\det(D_if_j) \notin P$ we deduce that $\det(D_ig_j) \notin P$, or in other words
$\rank J(g_1, \ldots, g_r; D_1, \ldots, D_r)(P) = r$. Therefore
\[
  \rank J(g_1, \ldots, g_s; D_1, \ldots, D_s)(P) = s.
\]
\end{proof}

If the above condition (1) holds, we say that the \emph{weak Jacobian condition} (WJ)
holds at $P$. If (WJ) holds at every $P \in \Spec R$ then we say that (WJ) holds in $R$.
If this holds, then for any $P$, $Q \in \Spec R$ with $Q \subset P$, setting $\htop Q = s$ we
have
\[
  R_P/QR_P \text{ is regular} \iff
  \text{there exist } D_1, \ldots, D_s \in \Der(R) \text{ and }
  f_1, \ldots, f_s \in Q \text{ such that } \det(D_if_j) \notin P,
\]
(the implication ($\Leftarrow$) is given by Theorem~\hyperref[thm:30.4]{4}). This statement is
the Jacobian criterion for regularity. We can use $\Der_k(R)$ in place of
$\Der(R)$, and we then write (WJ)$_k$.

\begin{theorem}\label{thm:30.7}
Let $(A, m)$ be an $n$-dimensional Noetherian local integral domain containing
$\mathbb{Q}$, and let $k \subset A$ be a subfield such that $\trdeg_k(A/m) = r < \infty$. Then
$\Der_k(A)$ is isomorphic to a submodule of $A^{n+r}$, and is therefore a finite
$A$-module, with
\[
  \rank \Der_k(A) \le \dim A + \trdeg_k(A/m).
\]
\end{theorem}

\begin{proof}
We write $A^*$ for the completion of $A$; choose a quasi-coefficient field $k'$ of $A$
containing $k$, and let $K$ be a coefficient field of $A^*$ containing $k'$. Let
$u_1, \ldots, u_r$ be a transcendence basis of $k'$ over $k$, and $x_1, \ldots, x_n$ a system
of parameters of $A$. We define $\varphi\colon \Der_k(A) \to A^{n+r}$ by
$\varphi(D) = (Du_1, \ldots, Du_r, Dx_1, \ldots, Dx_n)$; then $\varphi$ is $A$-linear, so that we
need only prove that it is injective. Suppose then that $Du_i = Dx_j = 0$ for all $i$
and $j$; then $D$ has a continuous extension to $A^*$, and this vanishes on
$B = K[[x_1, \ldots, x_n]]$. Now we do not know whether $A^*$ is an integral domain,
but it is finite as a $B$-module, so that any $a \in A$ is integral over $B$, and if
$f(X) \in B[X]$ has $a$ as a root, and has minimal degree, then $f(a) = 0$,
$f'(a) \ne 0$. Then $0 = D(f(a)) = f'(a) \cdot Da$, and $Da \in A$, so that, since a
non-zero element of $A$ cannot be a zero-divisor in $A^*$, we must have $Da = 0$.
Hence $D = 0$.
\end{proof}

\begin{remark}
If $k$ is an imperfect field then there are counter-examples even if
$k = A/m$: suppose $\chark k = p$ and $a \in k - k^p$, and set
$A = k[X, Y]_{(X, Y)}/(X^p + aY^p)$; then $\dim A = 1$ but
$\rank \Der_k(A) = 2$.
\end{remark}

\begin{theorem}\label{thm:30.8}
Let $(R, m)$ be a regular local ring containing $\mathbb{Q}$ and $k$ a quasi-coefficient
field of $R$. Then the following three conditions are equivalent:
\begin{enumerate}[(1)]
  \item (WJ)$_k$ holds at $m$;
  \item $\rank \Der_k(R) = \dim R$;
  \item (WJ)$_k$ holds at every $P \in \Spec R$.
\end{enumerate}
Furthermore, if these conditions hold, then for any $P \in \Spec R$, every element
of $\Der_k(R/P)$ is induced by an element of $\Der_k(R)$ and
\[
  \rank \Der_k(R/P) = \dim R/P.
\]
\end{theorem}

\begin{proof}
(1) $\Leftrightarrow$ (2) is known from Theorem~\hyperref[thm:30.6]{6}, and (1) is contained in
(3).

(1) $\Rightarrow$ (3) Write $R^*$ for the completion of $R$, and let $K$ be the coefficient
field of $R^*$ containing $k$; then by Theorem~\hyperref[thm:30.6]{6}, if
$x_1, \ldots, x_n$ is a regular system of parameters of $R$ then the derivations
$\partial/\partial x_i$ of $R^* = K[[x_1, \ldots, x_n]]$ belong to $\Der_k(R)$ for $1 \le i \le n$, and
form a basis of it. Now let $P \in \Spec R$, and write $\sigma\colon R \to R/P$ for the
natural homomorphism; to say that $D' \in \Der_k(R/P)$ is induced by
$D \in \Der_k(R)$ means that there is a commutative diagram
\[
\begin{tikzcd}
R \arrow[r, "\sigma"] \arrow[d, "D"'] & R/P \arrow[d, "D'"] \\
R \arrow[r, "\sigma"'] & R/P
\end{tikzcd}
\]
Suppose then that $D'$ is given; then $D'\sigma \in \Der_k(R, R/P)$, and this has a
unique extension to an element of $\Der_k(R^*, R^*/PR^*)$, so that it is uniquely
determined by its values on $x_1, \ldots, x_n$. Therefore if we choose
$b_1, \ldots, b_n \in R$ such that $D'(\sigma(x_i)) = \sigma(b_i)$, and set
$D = \sum b_i \partial/\partial x_i$, then $D'$ is induced by $D$.

Now $\Der_k(R, R/P)$ is a free $R/P$-module with basis $\bar{\sigma} \circ \partial/\partial x_i$ for
$1 \le i \le n$, and $\Der_k(R/P)$ can be identified with the submodule
\[
  N = \{D \in \Der_k(R, R/P) \mid D(f) = 0 \text{ for all } f \in P\}.
\]
Therefore, if $P = (f_1, \ldots, f_t)$ and $\htop P = r$ then
\[
  \rank \Der_k(R/P) = n - \rank J(f_1, \ldots, f_t; \partial/\partial x_1, \ldots,
      \partial/\partial x_n)(P);
\]
according to Theorem~\hyperref[thm:30.4]{4}, the right-hand side is $\ge n - r$, and by
Theorem~\hyperref[thm:30.7]{7} the left-hand side is $\le \dim R/P = n - r$, so we see that
\[
  \rank J(f_1, \ldots, f_t; \partial/\partial x_1, \ldots, \partial/\partial x_n)(P) = r,
\]
\[
  \rank \Der_k(R/P) = \dim R/P.
\]
\end{proof}

This theorem has many applications. For example, in the ring $R$ of convergent
power series in $n$ variables over $k = \mathbb{R}$ or $\mathbb{C}$ ($R$ is denoted by
$k\{X_1, \ldots, X_n\}$ in [N1], and elsewhere by $k\{X_1, \ldots, X_n\}$) we have
$\partial X_i/\partial X_j = \delta_{ij}$, so that the Jacobian criterion for regularity holds in
$R$. In characteristic $0$, the assumptions of the theorem are satisfied by the
formal power series ring over any field, or more generally by the formal power
series ring $R[[Y_1, \ldots, Y_n]]$ over a regular local ring $R$ which satisfies the
assumptions of the theorem. Even if $R$ is not local, but is a regular ring
containing a field $k$ of characteristic $0$, and such that the residue field at
every maximal ideal is algebraic over $k$, then if (WJ)$_k$ holds at every maximal
ideal of $R$, it in fact holds at every prime ideal of $R$, so that the Jacobian
criterion for regularity holds for example in rings such as
$k[X_1, \ldots, X_n] \otimes k[[Y_1, \ldots, Y_m]]$. A further extension of this theorem can be
found in Matsumura [2].

Now we are going to prove the results analogous to Theorem~\hyperref[thm:30.5]{5} for formal
power series rings over a field of characteristic $p$. This is a difficult theorem
obtained by Nagata [6]. First of all we have to do some preparatory work.

Let $k$ be a field and $k' \subset k$ a subfield; we say that $k'$ is \emph{cofinite} in $k$ if
$[k : k'] < \infty$. We say that a family $\mathcal{F} = \{k_\alpha\}_{\alpha \in I}$ of subfields of $k$ is
a \emph{directed family} if for any $\alpha, \beta \in I$ there exists $\gamma \in I$ such that
$k_\alpha, k_\beta \subset k_\gamma$.

Let $K$ be a field of characteristic $p$, and $k \subset K$ a subfield. Then there
exists a directed family $\mathcal{F} = \{k_\alpha\}_{\alpha \in I}$ of intermediate fields
$k \subset k_\alpha \subset K$ cofinite in $K$ such that $\bigcap_\alpha k_\alpha = k(K^p)$. To construct
this, we let $B$ be a fixed $p$-basis of $K$ over $k$, and let $I$ be the set of all
finite subsets of $B$; then we only need take $k_\alpha = k(K^p, B - \alpha)$ for
$\alpha \in I$.

\begin{lemma}\label{lem:30.3}
Let $K$ be a field, $(k_\alpha)_{\alpha \in I}$ a directed family of subfields of $K$, and set
$k = \bigcap_\alpha k_\alpha$. Then if $V$ is a vector space over $K$, and
$v_1, \ldots, v_n \in V$ are linearly independent over $k$, there exists $\alpha \in I$ such that
$v_1, \ldots, v_n$ are also linearly independent over $k_\alpha$.
\end{lemma}

\begin{proof}
For each $\alpha \in I$, write $q(\alpha)$ for the number of linearly independent
elements over $k_\alpha$ among $v_1, \ldots, v_n$; let $\alpha$ be such that $q(\alpha)$ is maximal,
and set $q = q(\alpha)$. Now if $q < n$, and we assume that $v_1, \ldots, v_q$ are linearly
independent over $k_\alpha$, we have $v_n = \sum_{i=1}^{q} c_iv_i$ with $c_i \in k_\alpha$. Since
the $v_i$ are linearly independent over $k$, at least one of the $c_i$ does not belong
to $k$, so that we can assume $c_1 \notin k$. Hence there exists $\beta \in I$ such that
$c_1 \notin k_\beta$, and also $\gamma \in I$ such that $k_\alpha, k_\beta \subset k_\gamma$, so that
$v_1, \ldots, v_q$ and $v_n$ are linearly independent over $k_\gamma$; this contradicts the
maximality of $q$, so that $q = n$.
\end{proof}

\begin{lemma}\label{lem:30.4}
Let $k \subset K$ be fields of characteristic $p$, and let $\mathcal{F} = \{k_\alpha\}_{\alpha \in I}$
be a directed family of intermediate fields $k \subset k_\alpha \subset K$; then the
following conditions are equivalent:
\begin{enumerate}[(1)]
  \item $\bigcap_\alpha k_\alpha(K^p) = k(K^p)$;
  \item the natural map $\Omega_{K/k} \to \varprojlim \Omega_{K/k_\alpha}$ is injective;
  \item if a finite subset $\{u_1, \ldots, u_n\} \subset K$ is $p$-independent over $k$, then
        it is also $p$-independent over $k_\alpha$ for some $\alpha$;
  \item there exists a $p$-basis $B$ of $K$ over $k$ such that every finite subset of
        $B$ is $p$-independent over $k_\alpha$ for some $\alpha$.
\end{enumerate}
\end{lemma}

\begin{proof}
(1) $\Rightarrow$ (3) If $\{u_1, \ldots, u_n\}$ is $p$-independent over $k$ then the $p^n$ monomials
$u_1^{i_1} \cdots u_n^{i_n}$ for $0 \le i_\nu < p$ are linearly independent over $k(K^p)$, so that by
the previous lemma they are also linearly independent over $k_\alpha(K^p)$ for some
$\alpha$.

(3) $\Rightarrow$ (4) is trivial; any $p$-basis will do.

(4) $\Rightarrow$ (2) If $0 \ne \omega \in \Omega_{K/k}$ and $B$ is a $p$-basis then there is a unique
expression $\omega = \sum c_id_{K/k}b_i$, with $\{b_1, \ldots, b_n\} \subset B$ a finite set and
$c_i \in K$. If we take $\alpha$ such that $b_1, \ldots, b_n$ are $p$-independent over
$k_\alpha$, then $d_{K/k_\alpha}b_i$ for $1 \le i \le n$ are linearly independent as elements of
$\Omega_{K/k_\alpha}$, so that $\omega$ has non-zero image in $\varprojlim \Omega_{K/k_\alpha}$.

(2) $\Rightarrow$ (1) Let $a \in K$ be such that $a \notin k(K^p)$; then $d_{K/k}a \ne 0$, so that
$d_{K/k}a \ne 0$ in $\Omega_{K/k_\alpha}$ for some $k_\alpha$, in other words $a \notin k_\alpha(K^p)$.
\end{proof}

\begin{lemma}\label{lem:30.5}
Let $k \subset K$ and $\mathcal{F} = \{k_\alpha\}_{\alpha \in I}$ be as in the previous lemma, and
assume that $\bigcap_\alpha k_\alpha(K^p) = k(K^p)$. Then if $L$ is an extension field of $K$
which is either separable over $K$ or finitely generated over $K$, we again have
$\bigcap_\alpha k_\alpha(L^p) = k(L^p)$.
\end{lemma}

\begin{proof}
(i) If $L/K$ is separable, choose a $p$-basis $B$ of $K/k$ and a $p$-basis $C$ of
$L/K$; then in view of the exact sequence (Theorem~\hyperref[thm:25.1]{25.1})
\[
  0 \to \Omega_{K/k} \otimes_K L \to \Omega_{L/k} \to \Omega_{L/K} \to 0,
\]
$B \cup C$ is a $p$-basis of $L/k$. If $b_1, \ldots, b_m \in B$ and $c_1, \ldots, c_n \in C$ are
finitely many distinct elements, then by the previous lemma, $b_1, \ldots, b_m$ are
$p$-independent over some $k_\alpha$, and from this (replacing $k$ by $k_\alpha$ in the
above exact sequence) we see that $b_1, \ldots, b_m, c_1, \ldots, c_n$ are
$p$-independent over $k_\alpha$.

(ii) If $L/K$ is finitely generated, then since any finitely generated extension
can be obtained as a succession of elementary extensions of type
\begin{enumerate}[(a)]
  \item $L = K(x)$ with $x$ separable over $K$,
  or
  \item $L = K(x)$ with $x^p = a \in K$,
\end{enumerate}
we need only consider case (b). We further divide this into two subcases: in
the first, $d_{K/k}a = 0$, and then from Theorem~\hyperref[thm:25.2]{25.2} and the fact that
$L \cong K[X]/(X^p - a)$ we get $\Omega_{L/k} = (\Omega_{K/k} \otimes L) \oplus Ldx$; from this
one sees easily that $\Omega_{L/k} \to \varprojlim \Omega_{L/k_\alpha}$ is injective. In the second
subcase, $d_{K/k}a \ne 0$ and now we have
$\Omega_{L/k} \cong ((\Omega_{K/k} \otimes L)/L \cdot d_{K/k}a) \oplus Ldx$. Hence if we take a
$p$-basis of $K/k$ of the form $\{a\} \cup B'$ with $a \notin B'$, then $\{x\} \cup B'$ is a
$p$-basis of $L/k$. Now $L/k$ satisfies condition (4) of the previous lemma, since
for $b_1, \ldots, b_m \in B'$, if we choose $\alpha$ such that $\{a, b_1, \ldots, b_m\} \subset K$
is $p$-independent over $k_\alpha$, then $\{x, b_1, \ldots, b_m\} \subset L$ is $p$-independent
over $k_\alpha$.
\end{proof}

\begin{lemma}\label{lem:30.6}
Let $K$ be a field of characteristic $p$, and let $\{K_\alpha\}$ be a directed family of
cofinite subfields $K_\alpha \subset K$ such that $\bigcap_\alpha K_\alpha = K^p$. Then if $L$ is a
finite field extension of $K$, there exists $\alpha$ such that
\[
  \rank_L \Omega_{L/K_\alpha} = \rank_L \Omega_{L/K}
\]
for all subfields $K' \subset K_\alpha$ with $[K : K'] < \infty$.
\end{lemma}

\begin{proof}
Let $K = K_0 \subset K_1 \subset \cdots \subset K_t = L$ be a chain of intermediate
fields such that $K_i = K_{i-1}(x_i)$, with $x_i$ either separable algebraic over
$K_{i-1}$ or $x_i^p \in K_{i-1}$. Then by the previous lemma we have
$\bigcap_\alpha K_\alpha(L^p) = L^p$, so that we need only prove the lemma for $t = 1$;
hence suppose that $L = K(x)$. If $L$ is separable over $K$ then in view of
$\Omega_{L/K'} = \Omega_{K/K'} \otimes_K L$, the assertion is clear (any $\alpha$ will do). Thus
suppose that $x^p = a \in K$, but $a \notin K^p$; then there exists $\alpha$ such that
$a \notin K_\alpha$. If $K' \subset K_\alpha$, then computing by means of
$L = K[X]/(X^p - a)$ we see that
\[
  \Omega_{L/K'} = ((\Omega_{K/K'} \otimes_K L) \oplus Ldx)/L \cdot d_{K/K'}a,
\]
and if $\rank \Omega_{L/K'} < \infty$ then $\rank \Omega_{L/K'} = \rank \Omega_{L/K}$.
\end{proof}

\begin{theorem}[\textbf{Nagata}]\label{thm:30.9}
Let $k$ be a field of characteristic $p$, $S = k[Y_1, \ldots, Y_m]$ and $P \in \Spec S$;
suppose that $\mathcal{F} = \{k_\alpha\}_{\alpha \in I}$ is a directed family of cofinite subfields of
$k$ such that $\bigcap_\alpha k_\alpha = k^p$. Then there exists $\alpha \in I$ such that for
every intermediate field $k^p \subset k' \subset k_\alpha$ with $[k : k'] < \infty$ the
following formula holds:
\[
  \rank \Der_{k'}(S/P) = \dim(S/P) + \rank \Der_{k'}(k).
\]
\end{theorem}

\begin{proof}
Set $A = S/P$, let $L$ be the field of fractions of $A$, and $\dim A = n$. Choose a
system of parameters $x_1, \ldots, x_n$ of $A$, set $B = k[x_1, \ldots, x_n]$, let $K$ be the
field of fractions of $B$ and $m_B$ its maximal ideal. Then $A$ is a finite $B$-module,
and hence $[L : K] < \infty$. If $k'$ is an intermediate field $k^p \subset k' \subset k$
with $[k : k'] = p^\rho < \infty$, and if $u_1, \ldots, u_\rho$ is a $p$-basis of $k$ over $k'$,
then $\{u_1, \ldots, u_\rho, x_1, \ldots, x_n\}$ is a $p$-basis of $B$ over
$C' = k'[x_1^p, \ldots, x_n^p]$, in the sense that $B$ is the free $C'$-module with basis
the set of $p$-monomials in $u_1, \ldots, u_\rho, x_1, \ldots, x_n$. A derivation from $B$ to
$B$ is continuous in the $m_B$-adic topology, and any element of $\Der_{k'}(B)$ is $0$ on
$C'$. Therefore $\Der_{k'}(B) = \Der_{k'}(B) = \Hom_B(\Omega_{B/C'}, B)$, and $\Omega_{B/C'}$ is
the free $B$-module of rank $n + \rho$ with basis $du_1, \ldots, du_\rho, dx_1, \ldots, dx_n$.
Therefore $\Der_{k'}(B)$ is also a free $B$-module of rank $n + \rho$, and the theorem
holds in case $A = B$.

We write $F'$ for the field of fractions of $C'$, and for $k_\alpha \in \mathcal{F}$ we set
$C_\alpha = k_\alpha[x_1^p, \ldots, x_n^p]$ and write $K_\alpha$ for the field of fractions of
$C_\alpha$. As above we have
\[
  \Der_{k'}(A) = \Der_{k'}(A) = \Hom_A(\Omega_{A/C'}, A),
\]
and since $A$ is a finite $C'$-module, $\Omega_{A/C'}$ is a finite $A$-module. Hence
\[
  \Der_{k'}(A) \otimes_A L = \Hom_L(\Omega_{A/C'} \otimes_A L, L) = \Hom_L(\Omega_{L/F'}, L),
\]
so that $\rank_L \Der_{k'}(A) = \rank_L \Omega_{L/F'}$. Moreover,
\[
  n + \rank \Der_{k'}(k) = \rank \Omega_{L/F'} = \rank \Omega_{K/F'},
\]
so that the conclusion of the theorem can be rewritten
\[
  \rank \Omega_{L/F'} = \rank \Omega_{K/F'}.
\]
Any element of $K_\alpha$ can be written with its denominator in
$k^p[x_1^p, \ldots, x_n^p] = B^p$. However, $B$ is faithfully flat over $C_\alpha$, so that
$K_\alpha \cap B = C_\alpha$. From this one deduces easily that
$\bigcap_\alpha K_\alpha = K^p$. Thus the theorem follows from the previous lemma.
\end{proof}

\begin{theorem}[\textbf{Nagata's Jacobian criterion}]\label{thm:30.10}
Let $k$ be a field of characteristic $p$, $S = k[[X_1, \ldots, X_n]]$, and let $I$, $P$ be
ideals of $S$ such that $I \subset P \in \Spec S$; set $S_P = R$, $R/IR = A$, and suppose
that $\htop IR = r$ and $I = (f_1, \ldots, f_t)$. Choose a $p$-basis $\{u_\gamma\}_{\gamma \in \Gamma}$ of
$k$, and define $D_\gamma \in \Der(k)$ by $D_\gamma(u_{\gamma'}) = \delta_{\gamma\gamma'}$. We make
any element $D \in \Der(k)$ act on the coefficients of power series, thus
extending $D$ to an element of $\Der(S)$.

Then the following conditions are equivalent:
\begin{enumerate}[(1)]
  \item $A$ is a regular local ring;
  \item there exists a finite number of elements $\alpha, \beta, \ldots, \gamma \in \Gamma$
        such that
        \[
          \rank J(f_1, \ldots, f_t; D_\alpha, \ldots, D_\gamma,
              \partial/\partial X_1, \ldots, \partial/\partial X_n)(P) = r.
        \]
\end{enumerate}
\end{theorem}

\begin{proof}
(2) $\Rightarrow$ (1) follows from Theorem~\hyperref[thm:30.4]{4}.

(1) $\Rightarrow$ (2) As we see from the proof of Lemma~\hyperref[lem:30.2]{2}, we need only
prove (2) in the case $I = P$. We let $\mathcal{F}$ be the family of subfields of $k$
obtained by adjoining all but a finite number of $\{u_\gamma\}_{\gamma \in \Gamma}$ to $k^p$;
then the conditions of Theorem~\hyperref[thm:30.9]{9} are satisfied, and there exists
$k' \in \mathcal{F}$ such that
\[
  \rank \Der_{k'}(S/P) = n - \htop P + \rank \Der_{k'}(k).
\]
If we set $[k : k'] = p^s$ then by construction there exist
$\gamma_1, \ldots, \gamma_s \in \Gamma$ such that $k = k'(u_{\gamma_1}, \ldots, u_{\gamma_s})$; now set
$C = k'[[X_1, \ldots, X_n]]$, so that $\Omega_{S/C}$ is the free $S$-module with basis
$du_{\gamma_1}, \ldots, du_{\gamma_s}, dX_1, \ldots, dX_n$, and arguing as in the proof of
Theorem~\hyperref[thm:30.8]{8}, we see that the derivations of $S/P$ over $k'$ are all induced by
derivations of $S$ over $k'$, and that if $g_1, \ldots, g_m$ are generators of $P$, we have
\[
  \rank J(g_1, \ldots, g_m; D_{\gamma_1}, \ldots, D_{\gamma_s},
      \partial/\partial X_1, \ldots, \partial/\partial X_n)(P) = \htop P.
\]
\end{proof}

\begin{remark}
Conditions (1) and (2) in Theorem~\hyperref[thm:30.10]{10} are of the same form as the
corresponding conditions in Theorem~\hyperref[thm:30.5]{5}. Condition (3) of
Theorem~\hyperref[thm:30.5]{5} is not applicable as it stands to the present situation, since
$\Omega_{A/k}$ is in general not a finite $A$-module. However, in general for a module
$M$ over a local ring $(R, m)$, if we write (in temporary notation)
$\tilde{M} = M/\bigcap_\nu m^\nu M$ for the associated separated module, then for any
$R$-module $N$ which is separated ($N = \tilde{N}$) we have
$\Hom_R(\tilde{M}, N) = \Hom_R(M, N)$. Therefore in the situation of the above theorem
we have $\Der_k(S) = \Hom_S(\tilde{\Omega}_{S/k}, S)$ and $\Der_k(A) = \Hom_A(\tilde{\Omega}_{A/k}, A)$,
and moreover $\tilde{\Omega}_{S/k}$ is a free $S$-module of rank $\rho + \rank \Omega_{k/\pi}$.
From this, using the same argument as in the proof of Theorem~\hyperref[thm:30.5]{5}, we see
that replacing $\Omega_{A/k}$ by $\tilde{\Omega}_{A/k}$ in (3) of Theorem~\hyperref[thm:30.5]{5}, this
condition is equivalent to (1) and (2) of Theorem~\hyperref[thm:30.10]{10}. Verifying this is a
suitable exercise for the reader.
\end{remark}

\begin{corollary*}
Let $A$ be a complete Noetherian local ring; then $\Reg(A)$ is an open subset of
$\Spec A$.
\end{corollary*}

\begin{proof}
If $A$ is equicharacteristic, and $k$ is a coefficient field of $A$, then $A$ is of the
form $A = S/I$ with $S = k[[X_1, \ldots, X_n]]$, so that by Theorems~\hyperref[thm:30.8]{8} and
\hyperref[thm:30.10]{10}, we see as in the corollary of Theorem~\hyperref[thm:30.5]{5} that $\Reg(A)$ is open.

If $A$ is of unequal characteristic, then by Theorem~\hyperref[thm:24.4]{24.4}, it is enough to
prove, under the assumption that $A$ is an integral domain, that $\Reg(A)$
contains a non-empty open subset of $\Spec A$. Now by
Theorem~\hyperref[thm:29.5]{29.5}, $A$ contains a regular local ring $B$, and is a finite module
over $B$, so that if we let $L$ and $K$ be the fields of fractions of $A$ and $B$,
then $L$ is a finite extension of $K$, and is separable since $\chark K = 0$. Replacing
$A$ by $A_b$, and $B$ by $B_b$ for some suitable $0 \ne b \in B$ we can assume that $A$ is
a free $B$-module (although $A$ and $B$ are no longer local rings, $B$ remains
regular). Suppose that $\omega_1, \ldots, \omega_m$ are a basis of $A$ as a $B$-module, and
consider the discriminant
\[
  d = \det(\tr_{L/K}(\omega_i\omega_j)).
\]
Let us prove that if $P \in \Spec A$ is such that $d \notin P$ then $P \in \Reg(A)$. Set
$p = P \cap B$; then $A_p$ is flat over $B_p$, and $B_p$ is regular, so that we need only
prove that the fibre $A_p \otimes_B \kappa(p)$ is regular. Now $A = \sum B\omega_i$, so
that $A \otimes_B \kappa(p) = \sum \kappa(p)\bar{\omega}_i$, and in $\kappa(p)$ we have
$\det(\tr(\bar{\omega}_i\bar{\omega}_j)) = \bar{d} \ne 0$; hence $A \otimes_B \kappa(p)$ is reduced,
and is therefore a direct product of fields. Therefore $A_p \otimes_B \kappa(p)$ is a
field; this proves that $\Reg(A)$ contains a non-empty open subset of $\Spec A$.
\end{proof}

% ============ 习题 ============
\section*{Exercises to \S\ref{sec:30}}

\begin{exercise}\label{ex:30.1-ex}
Let $(A, m)$ be a complete Noetherian local ring, and $D = (D_0, D_1, \ldots) \in
\mathrm{HS}(A)$. Suppose that $x \in m$ satisfies $D_1x = 1$ and $D_ix = 0$ for $i > 0$,
that is $E_t(x) = x + t$, and define $\varphi = E_{-t}$ by
$\varphi(a) = \sum_i (-x)^iD_ia$. Then $\varphi$ is an endomorphism of $A$ to $A$, with
$\Ker \varphi = xA$; and if we set $C = \Im \varphi$ then $A = C[x] \cong C[[X]]$.
\end{exercise}

\begin{exercise}\label{ex:30.2-ex}
If the conditions of Theorem~\hyperref[thm:30.5]{5} hold, is it true that $A$ is 0-smooth over
the $k'$ appearing in (3)?
\end{exercise}

\begin{exercise}\label{ex:30.3-ex}
Are conditions (1) and (3) of Theorem~\hyperref[thm:30.3]{3} still equivalent if
$S = k[[X_1, \ldots, X_n]]$?
\end{exercise}

\begin{exercise}\label{ex:30.4-ex}
Let $R$ be a regular ring containing a field of characteristic $0$; if (WJ) holds in
$R$ then it also holds in the polynomial ring $R[X_1, \ldots, X_n]$.
\end{exercise}
